\documentclass[../../main.tex]{subfiles}
% \documentclass[a4paper,10pt]{ltjsarticle}
\begin{document}
\rhead{東工大後期$1990$年$1$番}
\setcounter{equation}{0}
\setcounter{figure}{0}
\setcounter{table}{0}
\begin{oframed}
  \documentclass[../../main.tex]{subfiles}
\begin{document}

$(x+1)(x-2)$ の小数第 $1$ 位を四捨五入したものが $1+5x$ と等しくなるような実数 $x$ を求めよ．

\end{document}

\end{oframed}
\setlength{\columnseprule}{0.4pt}
\begin{multicols}{2}
  {\bf [解答]}

  実数 $x$ に対して $f(x) = (x+1)(x-2)$, $g(x) = 1+5x$ とおく．題意から，
  \begin{align}\label{eq:1}
    g(x) \in \mathbb{Z}, \quad g(x) - \dfrac{1}{2} \le f(x) < g(x) + \dfrac{1}{2}
  \end{align}
  をみたす $x \in \mathbb{R}$ をもとめればよい．
  $g(x) \in \mathbb{Z}$ から $5x \in \mathbb{Z}$．つまり，$x = \dfrac{t}{5}$ ($t \in \mathbb{Z}$)
  とかける．\cref{eq:1}に代入して$t$の条件式を求めると，
  \begin{align}
    t + \dfrac{1}{2} & \le \left(\dfrac{t}{5}+1\right)\left(\dfrac{t}{5}-2\right) < t+\dfrac{3}{2} \nonumber \\
    t + \dfrac{1}{2} & \le \dfrac{1}{25}t^2  - \dfrac{1}{5}t - 2 < t+\dfrac{3}{2} \nonumber      \\
    \therefore 125  & \le 2t^2 - 60t < 175 \label{eq:2}
  \end{align}
  を得る．

  ここで，二次関数$y = p(t) = 2t^2 - 60t$ のグラフの概形は\cref{1990-1:fig:1}のようになっており，
  \begin{align}
    & p(-3) = 198 & p(-1) = 62  \\
    & p(31) = 62  & p(33) = 198
  \end{align}
  だから，\cref{eq:2}を満たすような$t \in \mathbb{Z}$は $t=-2, 32$ である．したがって求めるべき$x=t/5$は$x = \dfrac{32}{5}, -\dfrac{2}{5}$ である． $\cdots$(答)

  \begin{figure}[H]
    \centering
    % Define values for clarity
    \def\yBoundaryLower{125}
    \def\yBoundaryUpper{175}
    \def\yAtIntPoints{128} % y-value for t=-2 and t=32

    % % Calculate t-values for y=125
    \pgfmathsetmacro{\tLowerSmallMinus}{15 - 0.5 * sqrt(1150)} % 15 - (5/2)*sqrt(46) approx -1.95
    \pgfmathsetmacro{\tUpperSmallPlus}{15 + 0.5 * sqrt(1150)} % 15 + (5/2)*sqrt(46) approx 31.95

    % % Calculate t-values for y=175
    \pgfmathsetmacro{\tLowerLargeMinus}{15 - 0.5 * sqrt(1250)} % 15 - (25/2)*sqrt(2) approx -2.675
    \pgfmathsetmacro{\tUpperLargePlus}{15 + 0.5 * sqrt(1250)} % 15 + (25/2)*sqrt(2) approx 32.675

    \begin{tikzpicture}[
        xscale=0.15,
        yscale=0.012,
        >=stealth,
      ]
      % Define t-range for plotting
      \def\tmin{-5}
      \def\tmax{36}
      \def\ymin{-20}
      \def\ymax{250}

      % Axes
      \draw[->, thick, gray!80!black] (\tmin, 0) -- (\tmax+2, 0) node[right] {$t$};
      \draw[->, thick, gray!80!black] (0, \ymin) -- (0, \ymax) node[above] {$y$};

      % Label the origin
      \node at (0,0) [below right=1pt, font=\small] {$O$};

      % Boundary dashed lines for y=125 and y=175
      \draw[dashed, gray!70!black] (\tmin, \yBoundaryLower) node[left, font=\scriptsize] {$125$} -- (\tmax, \yBoundaryLower) node[right, font=\scriptsize] {$y=125$};
      \draw[dashed, gray!70!black] (\tmin, \yBoundaryUpper) node[left, font=\scriptsize] {$175$} -- (\tmax, \yBoundaryUpper) node[right, font=\scriptsize] {$y=175$};

      % Plot the function y = 2t^2 - 60t
      \draw[blue!85!black, thick, domain=-3.5:33.5, samples=120] plot (\x, {2*\x*\x - 60*\x});
      \node[blue!85!black, font=\small] at (15, 60) {$y = 2t^2 - 60t$};

      % Highlight the parts of the parabola where 125 <= y < 175
      \draw[red!85!black, line width=1.5pt, domain=\tLowerLargeMinus:\tLowerSmallMinus, samples=50] plot (\x, {2*\x*\x - 60*\x});
      \draw[red!85!black, line width=1.5pt, domain=\tUpperSmallPlus:\tUpperLargePlus, samples=50] plot (\x, {2*\x*\x - 60*\x});

      % Horizontal connection between integer points at y=128
      \draw[densely dotted, thick, red!70!black] (-2, \yAtIntPoints) -- (32, \yAtIntPoints);
      \node[above right=1pt, font=\scriptsize, red!85!black] at (0, \yAtIntPoints) {$128$};

      % Vertical drop lines to t-axis
      \draw[densely dotted, thick, gray!70!black] (-2, 0) -- (-2, \yAtIntPoints);
      \node at (-2, 0) [below left=1pt, font=\small] {$-2$};
      \draw[densely dotted, thick, gray!70!black] (32, 0) -- (32, \yAtIntPoints);
      \node at (32, 0) [below=1pt, font=\small] {$32$};

      % Mark integer points (-2, 128) and (32, 128)
      \node[circle, fill=red!85!black, inner sep=1.4pt] at (-2, \yAtIntPoints) {};
      \node[above left=2pt, font=\scriptsize, red!85!black] at (-2, \yAtIntPoints) {$(-2, 128)$};

      \node[circle, fill=red!85!black, inner sep=1.4pt] at (32, \yAtIntPoints) {};
      \node[above right=2pt, font=\scriptsize, red!85!black] at (32, \yAtIntPoints) {$(32, 128)$};
    \end{tikzpicture}
    \caption{二次関数 $y = 2t^2 - 60t$ のグラフ}
    \label{1990-1:fig:1}
  \end{figure}

  {\bf [解説]}
  二次関数の問題．条件を素直に式に落としていけば解ける比較的容易な問題である．
  二次関数 $(x+1)(x-2)$ と一次関数 $1+5x$ がほぼ等しくなるような条件なので，解はこれらの交点に近くなるだろうというのが予想できる．
  実際にこれを解いてみると
  \begin{align}
    (x+1)(x-2)   & = 1 + 5x \\
    x^2 - 6x - 3 & = 0   \\
    x = 3 \pm 2\sqrt{3}
  \end{align}
  であり，$x \approx 6.46$ と $x \approx -0.46$ が得られる．
  本問題の解答である$x = \dfrac{32}{5}, -\dfrac{2}{5}$ はこれらの値にほぼ等しく，検算として利用できるだろう．

  \newpage
\end{multicols}
\end{document}
